\emph{Setting.} Coordinates, the Hermitian form $G$, the generators $T_{j} = \tfrac{I_{4}}{2}\sigma_{j}$ (the unit on the rung field and its coefficients is the field theory's unit $I_{4}$, which commutes with the rung vectors; $\C = \mathrm{span}\{1, I_{4}\}$) and $Y = \tfrac{I_{4}}{2}\mathbf{1}$ (acting by $X$ on $\alpha$ and by $\bar{X}$ on $\beta$), the lift family and the computation of the physical directions (the $G$-orthogonal complement of the gauge orbit in the tangent space of $L$, the neutral line $dc_{\lambda}/d\lambda$ removed) are those of \claimref{synthesis.rung-field-d-term-action-and-tree-masses-on-the-lift-locus} and \claimref{synthesis.rung-field-on-the-lift-locus}. The Yukawa form $M_{x}(\psi, \varphi) = B(\psi, x\varphi)$ on the five-mode module and the species of its entries are those of \claimref{standard-model.yukawa-pattern-two-rung-module}: a rung creation vector $v_{a}$ ($Y = +\tfrac12$) has entries exactly on $(L, \nu^{c})$ and $(Q, u^{c})$, a rung annihilation vector $w_{a}$ ($Y = -\tfrac12$) exactly on $(L, e^{c})$ and $(Q, d^{c})$. The conjugation $v_{j}^{*} = w_{j}$, conjugate-linear, is that of \claimref{standard-model.imported-tower-commuting-copies-fock-module}, and the readings are those of \claimref{standard-model.one-vev-vector-relations-breaking-pattern}, taken sector by sector: down-type entries from $M_{c}$, up-type entries from $M_{c^{*}}$. This per-sector reading is not $\su(2)_{R}$ covariant: by (i) its kernel is the $\alpha$ half, which the $\su(2)_{R}$ raising generator carries into the $\beta$ half, and the kernel of an equivariant map is invariant. The covariant completion puts one family matrix $Y$ on $c$ and one $\tilde{Y}$ on $c^{*}$ as wholes, $M_{d} = \beta Y + \bar{\alpha}\tilde{Y}$, $M_{u} = \alpha Y + \bar{\beta}\tilde{Y}$, whose leading mass reading at $|\alpha| \ll |\beta|$ is the per-sector one (\claimref{synthesis.rung-field-coupling-to-its-conjugate-is-the-second-family-matrix}). Every clause below concerns the per-sector reading.

\emph{(i).} The down-type sectors $(Q, d^{c})$, $(L, e^{c})$ of $M_{c}$ receive only the $w_{a}$ terms of $c$, whose coefficients are $\beta_{a}$; the up-type sectors $(Q, u^{c})$, $(L, \nu^{c})$ of $M_{c^{*}}$ receive only the $v_{a}$ terms of $c^{*} = \sum_{a}(\bar{\beta}_{a}v_{a} + \bar{\alpha}_{a}w_{a})$, whose coefficients are $\bar{\beta}_{a}$. So both readings are functions of $\beta$ alone: zero for $\beta = 0$ and unchanged by $\alpha \mapsto u\alpha$. The check computes every entry on the $16 \times 16$ basis pairs for two pure-$\alpha$ vectors (zero), three generic Gaussian-rational vectors (nonzero, with $\sum|\text{entry}|^{2} = 8|\beta|^{2}$ in each sector) and two rational phases $u = \tfrac35 + \tfrac45I_{4}$, $\tfrac{5}{13} - \tfrac{12}{13}I_{4}$ (entries unchanged), and confirms that $M_{c}$ itself has nonzero up-type entries on pure $\alpha$: only the $c^{*}$ reading removes $\alpha$ from the up-type sector.

\emph{(ii).} At $c_{\lambda}$ the charged plane is $\{\alpha' = a\,e_{5},\ \beta' = (H^{2}/\lambda^{2})\bar{a}\,e_{5}\}$ (stored clause (3) of the D-term result; the check recomputes it from the tangent space and the gauge orbit and verifies the form of $\beta'$). By (i) the entries of a direction $d$ in either sector have $\sum|\cdot|^{2} = 8|\beta(d)|^{2}$, so the ratio in the statement is
\[
\frac{|\beta(d)|^{2}/\mathsf{N}(d)}{|\beta(c_{\lambda})|^{2}/\mathsf{N}(c_{\lambda})} = \frac{(H^{4}/\lambda^{4})|a|^{2}}{\tfrac12|a|^{2}(1 + H^{4}/\lambda^{4})}\cdot\frac{\tfrac12(H^{4}/\lambda^{2} + \lambda^{2})}{\lambda^{2}} = \frac{H^{4}}{\lambda^{4}}.
\]
Dividing by $\mathsf{N}$ normalises each direction for the Hermitian kinetic term, as in the stored mass computation. The check verifies the ratio exactly in both sectors for $(H^{2}, \lambda) = (1, 3), (1, \tfrac72), (2, 5), (\tfrac14, 9)$. In the two-doublet reading: $\Phi_{1} = \alpha$ transforms by $X$, and so does $\Phi_{2} = \bar{\beta}$, both with hypercharge $+\tfrac12$; $\Phi_{2}^{\dagger}\Phi_{1} = \beta^{\mathsf T}\alpha = Q(c)$; at $c_{\lambda}$, $|\Phi_{1}| = H^{2}/\lambda$ and $|\Phi_{2}| = \lambda$, so $\tan\beta_{2} = \lambda^{2}/H^{2} = e^{2\eta}$, every fermion entry comes from $\Phi_{2}$ (one doublet coupling to all fermions), and the charged direction $(\Phi_{1}', \Phi_{2}') = (a, (H^{2}/\lambda^{2})a)e_{5}$ is the orthogonal complement of the eaten direction, as in a two-doublet model with that angle; the ratio above is $\cot^{2}\beta_{2}$.

\emph{(iii).} The vertex of a $W_{j}$ with the neutral mode and a charged direction $d$ comes from the cross term of the Hermitian kinetic term, with coefficient proportional to $G(dc_{\lambda}/d\lambda, T_{j}d) = -G(T_{j}\,dc_{\lambda}/d\lambda, d)$, since $T_{j}$ is $G$-antisymmetric. For $j = 1, 2$, $T_{j}$ carries the slot-$4$ vector $dc_{\lambda}/d\lambda = (-H^{2}/\lambda^{2}, -1)e_{4}$ into slot $5$, where the four real directions are the two eaten ones $T_{1}c_{\lambda}$, $T_{2}c_{\lambda}$ and the two physical charged ones; the check verifies that the two planes are $G$-orthogonal and span the image ($|P_{\text{phys}}w|^{2} + |P_{\text{eaten}}w|^{2} = G(w, w)$), and that
\[
|P_{\text{phys}}w|^{2}(\lambda^{4} - H^{4})^{2} = 4H^{4}\lambda^{4}\,|P_{\text{eaten}}w|^{2}, \qquad w = T_{j}\frac{dc_{\lambda}}{d\lambda},
\]
in the four frames. Since $\sinh 2\eta = (\lambda^{4} - H^{4})/(2H^{2}\lambda^{2})$, the ratio is $1/\sinh^{2}2\eta$: in two-doublet language $\cos^{2}(\beta_{2} - \alpha_{2})/\sin^{2}(\beta_{2} - \alpha_{2})$ with the neutral mode as the light scalar.

\emph{(iv).} $\mathsf{N}$, $m$ and $G$ are sums of $|\alpha_{a}|^{2}$, $|\beta_{a}|^{2}$ and of $\bar{\alpha}_{a}\alpha'_{a}$, $\bar{\beta}_{a}\beta'_{a}$, kept by a common phase on $\alpha$; $u$ commutes with $X$ and so with $T_{j}$ and $Y$; the moment maps are $G(c, -I_{4}T c)$, kept; $|Q|$ is kept and $Q \mapsto uQ$. The check verifies each on three generic vectors and two phases, including $Q(uc) \ne Q(c)$, and that the locus $\{Q = 0\}$ is kept. The D-term action is a function of the $\mu$'s and of the level; every $W(\mathsf{N})$ is a function of $\mathsf{N}$; the kinetic term is built from $G$ and the generators; the Yukawa entries are (i). So the one term of the stored action that the phase does not keep is the constraint $Q(c) = -H^{2}$ that defines $L$. At $H = 0$ the vacuum is $\alpha = 0$, $\beta = (-\lambda, 0)$ on $\{Q = 0\}$, and the check finds the charged plane with $\beta' = 0$, $\alpha'_{5} \ne 0$ at $\lambda = 3, \tfrac72$.

\emph{(v).} $\mathsf{N}(c_{\lambda}) = \tfrac12(H^{4}/\lambda^{2} + \lambda^{2}) < \lambda^{2}$ for $\lambda^{2} > H^{2}$, so $\cot\beta_{2} = H^{2}/\lambda^{2} < H^{2}/\mathsf{N}^{*} = g^{2}H^{2}/(2M_{W}^{2})$ with $M_{W}^{2} = \tfrac12 g^{2}\mathsf{N}^{*}$ (stored gauge form). The check bounds $g^{2} = 4\pi\alpha/\sin^{2}\theta_{W}$ above with $\pi < \tfrac{22}{7}$ and the stored $\alpha^{-1}(M_{Z}) = 127.951$, $\sin^{2}\theta_{W} = 0.23122$, takes $H^{2} = \Lambda(\hbar c)^{2}/3$ from the stored $\Lambda = 1.09 \times 10^{-52}$ m$^{-2}$ and $\hbar c = 197.3269804$ MeV fm, and the stored $M_{W} = 80.3692$ GeV, and finds $g^{2}H^{2}/(2M_{W}^{2}) < 4.66 \times 10^{-89}$. Since $\lambda^{2} = 2\mathsf{N}^{*} - H^{4}/\lambda^{2} < 2\mathsf{N}^{*}$, also $\cot\beta_{2} > H^{2}/(2\mathsf{N}^{*}) = g^{2}H^{2}/(4M_{W}^{2})$, and the check, with $\tfrac{333}{106} < \pi < \tfrac{22}{7}$, puts this lower end in $(2.32, 2.33) \times 10^{-89}$; the two differ from $\cot\beta_{2}$ by relative orders $1$ and $H^{4}/\mathsf{N}^{*2}$, so $\cot\beta_{2} \approx 2.3 \times 10^{-89}$ (computed: $H^{2} \approx 1.41 \times 10^{-84}$ GeV$^{2}$, $g^{2} \approx 0.425$, $v = \sqrt{2\mathsf{N}^{*}} \approx 246.6$ GeV). Here $H^{2} = \Lambda/3$, so $H$ is $H_{0}\sqrt{\Omega_{\Lambda}}$, not $H_{0}$; the two differ by a factor of order one. For a charged scalar coupling to $t\bar{b}$ through one doublet with angle $\beta_{2}$, $\Gamma(t \to bH^{+})/\Gamma(t \to bW^{+}) = \cot^{2}\beta_{2}\,(1 - x)^{2}/((1 - y)^{2}(1 + 2y))$ with $m_{b}$ neglected (Fact \claimref{synthesis.imported-two-doublet-top-to-charged-higgs-width-ratio}, the standard tree formulas); the check finds $K < \tfrac12$ at both ends of the bracket $m_{\pm}^{2} \in [13702, 13803]$ GeV$^{2}$ of \claimref{synthesis.rung-field-charged-scalar-sum-rule} (computed: $K \approx 0.33$), so the ratio is below $10^{-176}$ (computed: about $2 \times 10^{-178}$).

\emph{Reading (argued, not part of the statement).} (a) The ATLAS search for $t \to bH^{+}$, $H^{+} \to \tau\nu$ with $140$ fb$^{-1}$ at $13$ TeV bounds $\mathcal{B}(t \to bH^{+})\mathcal{B}(H^{+} \to \tau\nu)$ between $0.27\%$ and $0.02\%$ for $80$ to $160$ GeV (Fact \claimref{synthesis.imported-atlas-charged-higgs-top-decay-tau-nu-limits}); under the per-sector reading the branching is below $10^{-176}$ and the search says nothing. (b) Under the per-sector reading the phase on $\alpha$ is broken only through $Q(c) = -H^{2}$, and the charged pair is the lightest state carrying $\alpha$-number (the neutral component of $\alpha$ is fixed by the constraint): each of its couplings to lighter states computed here carries $H^{2}/\lambda^{2}$ or $1/\sinh 2\eta$, and its vertices with $W^{\pm}Z$ and $W^{\pm}\gamma$ vanish at tree level (the check computes the seagull coefficients at $c_{\lambda}$). Its partial widths are then of order $10^{-176}$ times a weak width, so it is stable on every time scale that matters. (c) A stable charged scalar near $117$ GeV is pair-produced through $\gamma^{*}/Z^{*}$ at hadron colliders and is not reached at LEP ($2m_{\pm} > 209$ GeV). CMS excluded Drell-Yan pair-produced long-lived staus below $339$ GeV at $7$ and $8$ TeV, with an observed limit $0.16$ of the predicted cross-section at the lowest tabulated mass $126$ GeV (Fact \claimref{synthesis.imported-long-lived-charged-particle-stau-limits}(i)); ATLAS at $13$ TeV with $36.1$ fb$^{-1}$ set $430$ GeV for staus, on a grid starting at $287$ GeV (Fact \claimref{synthesis.imported-long-lived-charged-particle-stau-limits}(ii)). The charged pair's $Z$ coupling, $\tfrac12 - \sin^{2}\theta_{W}$, is comparable to a right-handed stau's, $\sin^{2}\theta_{W}$, and its mass is lower, so its cross-section is comparable or larger. (d) The two readings then give two branches. Under the per-sector reading the couplings are suppressed by $e^{-2\eta}$ and the pair is stable, for every $W(\mathsf{N})$ and every level; near $117$ GeV it is excluded by the searches for stable charged particles of (c), a reinterpretation of the stau limits and not a dedicated analysis, and a population of stable charged particles left by a standard hot early universe is also excluded by searches for anomalous heavy hydrogen (memory, flagged). Under the covariant completion the $\alpha$ half carries the same $Y$ and $\tilde{Y}$ that give the masses, so the pair couples with the strength of the mass couplings, of order $m_{t}/v$ to $t$ and $b$, and searches in top decays apply: with a coupling of the size of the top's mass coupling the width ratio is $K$, about $0.33$ at $117$ GeV, a branching $K/(1 + K)$ of about $25\%$ (computed; the check brackets it in $(0.24, 0.25)$ at both ends of the mass bracket), far above the limits of (a) and of the corresponding searches in $H^{+} \to c\bar{s}$ (memory, flagged), so the $117$ GeV point is excluded on that branch too.
